\clearpage
\SourceLeaf{174}\centerline{\Large BRIEFWECHSEL.}

\bigskip
\hrule
\bigskip

\centerline{[LEGENDRES THEORIE DER PARALLELEN.]}

\bigskip
\hrule
\bigskip

\centerline{[\textsc{Gerling} \emph{an} \textsc{Gauss}. \ \emph{Kassel, 11. März 1816.}]}

\smallskip

\{.\ .\ .\ .\ .\ .\ .\ Am Schluss dieses Briefes fällt mir ein, dass ich schon oft vergessen habe, Sie um Ihr Urteil über \textsc{Legendres} Theorie der Parallelen in seinen élémens de géom. zu bitten. --- Er definiert die gerade Linie als den kürzesten Weg zwischen zwei Punkten und beweist mithilfe dieser Definition, dass die Summe der 3 Winkel im Dreieck \textls{nicht größer} sein kann als $2R$. Danach beweist er, dass sie auch \textls{nicht kleiner} sein könne als $2R$, wobei aber vorausgesetzt wird, dass man eine Linie durch einen Punkt zwischen den Schenkeln eines Winkels, der $< \tfrac{2}{3}R$ ist, immer so legen könne, dass beide Schenkel geschnitten werden. Diese Voraussetzung rechtfertigt er in einer Anmerkung (pag. 280, 6te edit. 1806) durch das Einholen und Überholen des Punktes durch Verbindungslinien zwischen gleich weit vom Scheitel entfernten Punkten auf den Schenkeln. --- In diesem letzten Schritt scheint mir aber derselbe
Fehler zu stecken, dessen Sie mich bei meinem letzten Aufenthalt in Göttingen überführten. Er erklärt es für \textls{assez évident} und glaubt, man könne es nicht strenger begründen, ohne von einer anderen Definition der geraden Linie auszugehen. --- Danach aber zeigt er auf eine andere, mir neue und, wie mir scheint, stringente Weise, dass die Summe der 3 Winkel im Dreieck $= 2R$ sein müsse. Ungefähr so geht er vor. Durch 2 Winkel und die zwischenliegende Seite $A, B, c$ ist das ganze Dreieck bestimmt, also
\[
C = \varphi(A,B,c).
\]
